% Options for packages loaded elsewhere
\PassOptionsToPackage{unicode}{hyperref}
\PassOptionsToPackage{hyphens}{url}
\PassOptionsToPackage{dvipsnames,svgnames,x11names}{xcolor}
%
\documentclass[
  10pt,
  a4paper,
  DIV=11]{scrartcl}

\usepackage{amsmath,amssymb}
\usepackage{iftex}
\ifPDFTeX
  \usepackage[T1]{fontenc}
  \usepackage[utf8]{inputenc}
  \usepackage{textcomp} % provide euro and other symbols
\else % if luatex or xetex
  \usepackage{unicode-math}
  \defaultfontfeatures{Scale=MatchLowercase}
  \defaultfontfeatures[\rmfamily]{Ligatures=TeX,Scale=1}
\fi
\usepackage{lmodern}
\ifPDFTeX\else  
    % xetex/luatex font selection
  \setmainfont[]{Arial}
\fi
% Use upquote if available, for straight quotes in verbatim environments
\IfFileExists{upquote.sty}{\usepackage{upquote}}{}
\IfFileExists{microtype.sty}{% use microtype if available
  \usepackage[]{microtype}
  \UseMicrotypeSet[protrusion]{basicmath} % disable protrusion for tt fonts
}{}
\makeatletter
\@ifundefined{KOMAClassName}{% if non-KOMA class
  \IfFileExists{parskip.sty}{%
    \usepackage{parskip}
  }{% else
    \setlength{\parindent}{0pt}
    \setlength{\parskip}{6pt plus 2pt minus 1pt}}
}{% if KOMA class
  \KOMAoptions{parskip=half}}
\makeatother
\usepackage{xcolor}
\usepackage[top=30mm,left=20mm,heightrounded]{geometry}
\setlength{\emergencystretch}{3em} % prevent overfull lines
\setcounter{secnumdepth}{3}
% Make \paragraph and \subparagraph free-standing
\ifx\paragraph\undefined\else
  \let\oldparagraph\paragraph
  \renewcommand{\paragraph}[1]{\oldparagraph{#1}\mbox{}}
\fi
\ifx\subparagraph\undefined\else
  \let\oldsubparagraph\subparagraph
  \renewcommand{\subparagraph}[1]{\oldsubparagraph{#1}\mbox{}}
\fi

\usepackage{color}
\usepackage{fancyvrb}
\newcommand{\VerbBar}{|}
\newcommand{\VERB}{\Verb[commandchars=\\\{\}]}
\DefineVerbatimEnvironment{Highlighting}{Verbatim}{commandchars=\\\{\}}
% Add ',fontsize=\small' for more characters per line
\usepackage{framed}
\definecolor{shadecolor}{RGB}{241,243,245}
\newenvironment{Shaded}{\begin{snugshade}}{\end{snugshade}}
\newcommand{\AlertTok}[1]{\textcolor[rgb]{0.68,0.00,0.00}{#1}}
\newcommand{\AnnotationTok}[1]{\textcolor[rgb]{0.37,0.37,0.37}{#1}}
\newcommand{\AttributeTok}[1]{\textcolor[rgb]{0.40,0.45,0.13}{#1}}
\newcommand{\BaseNTok}[1]{\textcolor[rgb]{0.68,0.00,0.00}{#1}}
\newcommand{\BuiltInTok}[1]{\textcolor[rgb]{0.00,0.23,0.31}{#1}}
\newcommand{\CharTok}[1]{\textcolor[rgb]{0.13,0.47,0.30}{#1}}
\newcommand{\CommentTok}[1]{\textcolor[rgb]{0.37,0.37,0.37}{#1}}
\newcommand{\CommentVarTok}[1]{\textcolor[rgb]{0.37,0.37,0.37}{\textit{#1}}}
\newcommand{\ConstantTok}[1]{\textcolor[rgb]{0.56,0.35,0.01}{#1}}
\newcommand{\ControlFlowTok}[1]{\textcolor[rgb]{0.00,0.23,0.31}{#1}}
\newcommand{\DataTypeTok}[1]{\textcolor[rgb]{0.68,0.00,0.00}{#1}}
\newcommand{\DecValTok}[1]{\textcolor[rgb]{0.68,0.00,0.00}{#1}}
\newcommand{\DocumentationTok}[1]{\textcolor[rgb]{0.37,0.37,0.37}{\textit{#1}}}
\newcommand{\ErrorTok}[1]{\textcolor[rgb]{0.68,0.00,0.00}{#1}}
\newcommand{\ExtensionTok}[1]{\textcolor[rgb]{0.00,0.23,0.31}{#1}}
\newcommand{\FloatTok}[1]{\textcolor[rgb]{0.68,0.00,0.00}{#1}}
\newcommand{\FunctionTok}[1]{\textcolor[rgb]{0.28,0.35,0.67}{#1}}
\newcommand{\ImportTok}[1]{\textcolor[rgb]{0.00,0.46,0.62}{#1}}
\newcommand{\InformationTok}[1]{\textcolor[rgb]{0.37,0.37,0.37}{#1}}
\newcommand{\KeywordTok}[1]{\textcolor[rgb]{0.00,0.23,0.31}{#1}}
\newcommand{\NormalTok}[1]{\textcolor[rgb]{0.00,0.23,0.31}{#1}}
\newcommand{\OperatorTok}[1]{\textcolor[rgb]{0.37,0.37,0.37}{#1}}
\newcommand{\OtherTok}[1]{\textcolor[rgb]{0.00,0.23,0.31}{#1}}
\newcommand{\PreprocessorTok}[1]{\textcolor[rgb]{0.68,0.00,0.00}{#1}}
\newcommand{\RegionMarkerTok}[1]{\textcolor[rgb]{0.00,0.23,0.31}{#1}}
\newcommand{\SpecialCharTok}[1]{\textcolor[rgb]{0.37,0.37,0.37}{#1}}
\newcommand{\SpecialStringTok}[1]{\textcolor[rgb]{0.13,0.47,0.30}{#1}}
\newcommand{\StringTok}[1]{\textcolor[rgb]{0.13,0.47,0.30}{#1}}
\newcommand{\VariableTok}[1]{\textcolor[rgb]{0.07,0.07,0.07}{#1}}
\newcommand{\VerbatimStringTok}[1]{\textcolor[rgb]{0.13,0.47,0.30}{#1}}
\newcommand{\WarningTok}[1]{\textcolor[rgb]{0.37,0.37,0.37}{\textit{#1}}}

\providecommand{\tightlist}{%
  \setlength{\itemsep}{0pt}\setlength{\parskip}{0pt}}\usepackage{longtable,booktabs,array}
\usepackage{calc} % for calculating minipage widths
% Correct order of tables after \paragraph or \subparagraph
\usepackage{etoolbox}
\makeatletter
\patchcmd\longtable{\par}{\if@noskipsec\mbox{}\fi\par}{}{}
\makeatother
% Allow footnotes in longtable head/foot
\IfFileExists{footnotehyper.sty}{\usepackage{footnotehyper}}{\usepackage{footnote}}
\makesavenoteenv{longtable}
\usepackage{graphicx}
\makeatletter
\def\maxwidth{\ifdim\Gin@nat@width>\linewidth\linewidth\else\Gin@nat@width\fi}
\def\maxheight{\ifdim\Gin@nat@height>\textheight\textheight\else\Gin@nat@height\fi}
\makeatother
% Scale images if necessary, so that they will not overflow the page
% margins by default, and it is still possible to overwrite the defaults
% using explicit options in \includegraphics[width, height, ...]{}
\setkeys{Gin}{width=\maxwidth,height=\maxheight,keepaspectratio}
% Set default figure placement to htbp
\makeatletter
\def\fps@figure{htbp}
\makeatother

%% load packages
\usepackage{sectsty}
\usepackage{fontspec}
\usepackage{scrlayer-scrpage}
\usepackage{fontspec}

% Define title formatting

%% Set font sizes for sections and subsections
\sectionfont{\fontsize{11}{16}\selectfont\bfseries} % Arial 14pt bold
\subsectionfont{\fontsize{10}{14}\selectfont} % Arial 12pt

\lehead[test1]{\pagemark}
\cehead[]{}
\rehead[test 2]{\leftmark}
\lohead[{\fontsize{8}{10}\selectfont\upshape\textbf{Institut Sekundarstufe I}} \\
\fontsize{8}{10}\selectfont\upshape Fabrikstrasse 8, CH-3012 Bern \\
\fontsize{8}{10}\selectfont\upshape T +41 31 309 21 15, contactdesk@phbern.ch, www.phbern.ch
]{\rightmark}
\rohead[PHBern]{}
\KOMAoption{captions}{tableheading}
\makeatletter
\@ifpackageloaded{tcolorbox}{}{\usepackage[skins,breakable]{tcolorbox}}
\@ifpackageloaded{fontawesome5}{}{\usepackage{fontawesome5}}
\definecolor{quarto-callout-color}{HTML}{909090}
\definecolor{quarto-callout-note-color}{HTML}{0758E5}
\definecolor{quarto-callout-important-color}{HTML}{CC1914}
\definecolor{quarto-callout-warning-color}{HTML}{EB9113}
\definecolor{quarto-callout-tip-color}{HTML}{00A047}
\definecolor{quarto-callout-caution-color}{HTML}{FC5300}
\definecolor{quarto-callout-color-frame}{HTML}{acacac}
\definecolor{quarto-callout-note-color-frame}{HTML}{4582ec}
\definecolor{quarto-callout-important-color-frame}{HTML}{d9534f}
\definecolor{quarto-callout-warning-color-frame}{HTML}{f0ad4e}
\definecolor{quarto-callout-tip-color-frame}{HTML}{02b875}
\definecolor{quarto-callout-caution-color-frame}{HTML}{fd7e14}
\makeatother
\makeatletter
\@ifpackageloaded{caption}{}{\usepackage{caption}}
\AtBeginDocument{%
\ifdefined\contentsname
  \renewcommand*\contentsname{Inhaltsverzeichnis}
\else
  \newcommand\contentsname{Inhaltsverzeichnis}
\fi
\ifdefined\listfigurename
  \renewcommand*\listfigurename{Abbildungsverzeichnis}
\else
  \newcommand\listfigurename{Abbildungsverzeichnis}
\fi
\ifdefined\listtablename
  \renewcommand*\listtablename{Tabellenverzeichnis}
\else
  \newcommand\listtablename{Tabellenverzeichnis}
\fi
\ifdefined\figurename
  \renewcommand*\figurename{Abbildung}
\else
  \newcommand\figurename{Abbildung}
\fi
\ifdefined\tablename
  \renewcommand*\tablename{Tabelle}
\else
  \newcommand\tablename{Tabelle}
\fi
}
\@ifpackageloaded{float}{}{\usepackage{float}}
\floatstyle{ruled}
\@ifundefined{c@chapter}{\newfloat{codelisting}{h}{lop}}{\newfloat{codelisting}{h}{lop}[chapter]}
\floatname{codelisting}{Listing}
\newcommand*\listoflistings{\listof{codelisting}{Listingverzeichnis}}
\usepackage{amsthm}
\theoremstyle{definition}
\newtheorem{exercise}{Übungsaufgabe}[section]
\theoremstyle{remark}
\AtBeginDocument{\renewcommand*{\proofname}{Beweis}}
\newtheorem*{remark}{Anmerkung}
\newtheorem*{solution}{Lösung}
\newtheorem{refremark}{Anmerkung}[section]
\newtheorem{refsolution}{Lösung}[section]
\makeatother
\makeatletter
\makeatother
\makeatletter
\@ifpackageloaded{caption}{}{\usepackage{caption}}
\@ifpackageloaded{subcaption}{}{\usepackage{subcaption}}
\makeatother
\ifLuaTeX
\usepackage[bidi=basic]{babel}
\else
\usepackage[bidi=default]{babel}
\fi
\babelprovide[main,import]{ngerman}
\ifPDFTeX
\else
\babelfont{rm}[]{Arial}
\fi
% get rid of language-specific shorthands (see #6817):
\let\LanguageShortHands\languageshorthands
\def\languageshorthands#1{}
\ifLuaTeX
  \usepackage{selnolig}  % disable illegal ligatures
\fi
\usepackage{bookmark}

\IfFileExists{xurl.sty}{\usepackage{xurl}}{} % add URL line breaks if available
\urlstyle{same} % disable monospaced font for URLs
\hypersetup{
  pdftitle={Übungen mit Conditionals},
  pdfauthor={Richard Conrardy},
  pdflang={de},
  colorlinks=true,
  linkcolor={(3,1.5)},
  filecolor={Maroon},
  citecolor={Blue},
  urlcolor={Blue},
  pdfcreator={LaTeX via pandoc}}

\title{Übungen mit Conditionals}
\author{Richard Conrardy}
\date{18.07.2024}

\begin{document}
\maketitle

\begin{exercise}[Gerade oder ungerade
Zahl]\protect\hypertarget{exr-line}{}\label{exr-line}

Schreibe ein Programm, das eine ganze Zahl als Eingabe nimmt und
überprüft, ob die Zahl gerade oder ungerade ist.

\end{exercise}

\begin{exercise}[Positiv, negativ oder
null]\protect\hypertarget{exr-line}{}\label{exr-line}

Schreibe ein Programm, das eine ganze Zahl als Eingabe nimmt und
überprüft, ob die Zahl positiv, negativ oder null ist.

\end{exercise}

\begin{exercise}[Zahlenmengen]\protect\hypertarget{exr-line}{}\label{exr-line}

Schreiben Sie ein Programm, welches prüft, ob eine Zahl in
\(\mathbb{N}\), \(\mathbb{Z}\) oder \(\mathbb{Q}\) ist.

\end{exercise}

\begin{exercise}[Mehrfache von drei und
fünf]\protect\hypertarget{exr-line}{}\label{exr-line}

Schreibe ein Programm, das eine ganze Zahl als Eingabe nimmt und
überprüft, ob die Zahl sowohl ein Vielfaches von 3 als auch von 5 ist.

\end{exercise}

\begin{exercise}[Schaltjahr oder
nicht]\protect\hypertarget{exr-line}{}\label{exr-line}

Schreibe ein Programm, das ein Jahr als Eingabe nimmt und überprüft, ob
es sich um ein Schaltjahr handelt.

\begin{quote}
Die gregorianische Schalttagsregelung besteht aus folgenden drei
einzelnen Regeln: 1. Die durch 4 ganzzahlig teilbaren Jahre sind,
abgesehen von den folgenden Ausnahmen, Schaltjahre. 2. Säkularjahre,
also die Jahre, die ein Jahrhundert abschließen (z.B. 1800, 1900, 2100
und 2200), sind, abgesehen von der folgenden Ausnahme, keine
Schaltjahre. 3. Die durch 400 ganzzahlig teilbaren Säkularjahre, zum
Beispiel das Jahr 2000, sind jedoch Schaltjahre.

Quelle: https://de.wikipedia.org/wiki/Schaltjahr
\end{quote}

\end{exercise}

\begin{exercise}[Abweichung]\protect\hypertarget{exr-line}{}\label{exr-line}

Schreiben Sie ein Programm, welche zwei Zahlen nimmt und augibt um
wieviel Prozent die grössere Zahl grösser ist als die kleinere Zahl.

\end{exercise}

\begin{exercise}[Eintritspreis basierend auf dem
Alter]\protect\hypertarget{exr-line}{}\label{exr-line}

Schreiben Sie ein Programm, das das Alter einer Person als Eingabe nimmt
und den Eintrittspreis basierend auf dem Alter berechnet: unter 18:
5Chf, 18-64: 10Chf, 65 und älter: 7Chf.

\end{exercise}

\begin{exercise}[Maximalwert]\protect\hypertarget{exr-line}{}\label{exr-line}

Schreiben Sie ein Programm, das drei Zahlen als Eingabe nimmt und den
größten Wert ausgibt.

\begin{tcolorbox}[enhanced jigsaw, colback=white, opacitybacktitle=0.6, left=2mm, breakable, rightrule=.15mm, leftrule=.75mm, title=\textcolor{quarto-callout-tip-color}{\faLightbulb}\hspace{0.5em}{Maximalfunktion}, coltitle=black, toptitle=1mm, colbacktitle=quarto-callout-tip-color!10!white, titlerule=0mm, bottomrule=.15mm, bottomtitle=1mm, colframe=quarto-callout-tip-color-frame, toprule=.15mm, opacityback=0, arc=.35mm]

Die Maximalfunktion gibt sofort den grössten Wert aus: \texttt{max()}

\end{tcolorbox}

\end{exercise}

\begin{exercise}[Berechnung der
Quadratwurzel]\protect\hypertarget{exr-line}{}\label{exr-line}

Schreiben Sie ein Programm, das eine Zahl als Eingabe nimmt und die
Quadratwurzel dieser Zahl berechnet. Überprüfen Sie dabei, ob die Zahl
positiv ist.

\end{exercise}

\begin{exercise}[Dreieckstyp]\protect\hypertarget{exr-line}{}\label{exr-line}

Schreiben Sie ein Programm, das die Längen der drei Seiten eines
Dreiecks als Eingabe nimmt und überprüft, ob es sich um ein Dreieck
handeln kann und falls ja, ob es sich um ein gleichseitiges,
gleichschenkliges oder unregelmäßiges Dreieck handelt.

\end{exercise}

\begin{exercise}[Addition mit
Überprüfung]\protect\hypertarget{exr-line}{}\label{exr-line}

Schreibe ein Programm, das zwei Zahlen als Eingabe nimmt, ihre Summe
berechnet und überprüft, ob das vom Benutzer eingegebene Ergebnis
korrekt ist.

\end{exercise}

\begin{exercise}[Zufällige Addition mit
Überprüfung]\protect\hypertarget{exr-line}{}\label{exr-line}

Schreibe ein Programm, das zwei zufällige Zahlen zwischen 1 und 100
generiert, ihre Summe berechnet und überprüft, ob das vom Benutzer
eingegebene Ergebnis korrekt ist.

\begin{tcolorbox}[enhanced jigsaw, colback=white, opacitybacktitle=0.6, left=2mm, breakable, rightrule=.15mm, leftrule=.75mm, title=\textcolor{quarto-callout-tip-color}{\faLightbulb}\hspace{0.5em}{Random-Package}, coltitle=black, toptitle=1mm, colbacktitle=quarto-callout-tip-color!10!white, titlerule=0mm, bottomrule=.15mm, bottomtitle=1mm, colframe=quarto-callout-tip-color-frame, toprule=.15mm, opacityback=0, arc=.35mm]

Das Python-Package random wird genutzt um (pseudo-)zufällige Zahlen
auszugeben. Jedes Mal wenn das Programm laufen gelassen werden neue
zufällige Zahlen mit Hilfe der aktuellen Uhrzeit produziert.

\begin{Shaded}
\begin{Highlighting}[]
\ImportTok{import}\NormalTok{ random}
\NormalTok{zahl1 }\OperatorTok{=}\NormalTok{ random.randint(}\DecValTok{1}\NormalTok{, }\DecValTok{100}\NormalTok{)}
\NormalTok{zahl2 }\OperatorTok{=}\NormalTok{ random.random()}
\NormalTok{zahl3 }\OperatorTok{=}\NormalTok{ random.choice(}\OperatorTok{{-}}\DecValTok{1}\NormalTok{,}\DecValTok{1}\NormalTok{)}
\BuiltInTok{print}\NormalTok{(zahl1,zahl2,zahl3)}
\end{Highlighting}
\end{Shaded}

\end{tcolorbox}

\end{exercise}

\begin{exercise}[Zufällige Multiplikation mit
Bewertung]\protect\hypertarget{exr-line}{}\label{exr-line}

Schreiben Sie ein Programm, das zwei zufällige Zahlen zwischen 1 und 100
generiert, ihr Produkt berechnet und überprüft, um wieviel das vom
Benutzer eingegebene Resultat vom echten Resultat abweicht. Bei einer
Eingabe von weniger als 2\% soll ``sehr gut'' ausgegeben werden, bei
einer Abweichung ovn weniger als 7\% soll ``gut'' augegeben werden und
bei einer Abweichung von mehr als 20\% soll ``schwach'' ausgegeben
werden. In allen anderen Fällen soll ``gut genug'' ausgegeben werden.

\end{exercise}



\end{document}
